\scanpage{56}

\margin{\fbox{\parbox{.42\linewidth}{on suppose \unsure{ici} que les bitors. birigidifiés sur $A\times A$ sont des biextensions.}}}
Reprenons les not. \unsure{et A} de la fin de b), On va définir un hom en sens inverse
\diagram{\begin{tikzcd}[column sep=large]
R^1f_*(C_A) & \underline{\operatorname{Tors}}\,\underline{\mathrm{birig}}\,\underline{\mathrm{sym}}(A,A;C)\arrow[l,bend right=15,"\Delta^*"']
\end{tikzcd}}
\margin{\unsure{inutile} pour définir $\Delta^*$}
en utilisant l'application diagonale, $\Delta:A\times_S A\mathbin{-}A$.

Calculons le composé \inserted{$h\Delta^*j'$} de $\Delta^*j'$ avec $h:R^1f_*(C_A)\mathbin{-}\underline{\operatorname{NS}}(A)$, on aura d'abord
\[
\Delta^*j(L)\simeq(2\operatorname{id}_A)^*(L)L^{-2}
\]
\unsure{Je} dis que c'est
\[
\equiv L^2\quad\{\mathrm{mod}\ A'=\underline{\operatorname{Ext}}^1(A,C)\}.
\]
Cela résulte du fait suivant

\uline{Lemme} Pour tout entier $n$, $(n\operatorname{id}_A)^*(L)\equiv n^2L$,
\deleted{\unsure{pour L dans}} plus généralement, si $A,B$ sont comme \unsure{énoncés}, alors
\[
\operatorname{Hom}(A,B)\longrightarrow\operatorname{Hom}(\underline{\operatorname{NS}}(B),\underline{\operatorname{NS}}(A))
\]
est une application \uline{quadratique} i.e.
\[
\varphi(u+v)-\varphi(u-v)=2(\varphi(u)+\varphi(v)).
\]
\deleted{\unsure{En ef}} \uline{Dém} En effet, cette application se factorise en
\diagram{\begin{tikzcd}[column sep=small]
\operatorname{Hom}(A,B)\arrow[r,dash] & \operatorname{Hom}(\underline{\operatorname{Bitors}}(B,B;C),\underline{\operatorname{Bitors}}(A,A;C))^!\arrow[d,dash,"\text{linéaire}"]\\
 & \operatorname{Hom}(\underline{\operatorname{NS}}(B),\underline{\operatorname{NS}}(A))
\end{tikzcd}}
où l'exposant ! indique les hom qui envoient $\underline{\operatorname{NS}}(B)$ dans $\underline{\operatorname{NS}}(A)$. \deleted{\unsure{On trouve}} Or il est évident que la première flèche est quadratique.

\margin{55}
